\newpage
\judulword{11}{0}{E. USULAN PENGEMBANGAN DAN KESIMPULAN}
\judulword{10}{4}{Usulan Pengembangan Penelitian}
\pertanyaan{9}{1. Judul/tema pengembangan yang diusulkan:}
\teksword{8.5}{0}{.........}
\pertanyaan{9}{2. Permasalahan yang akan diselesaikan:}
\teksword{8.5}{0}{.........}
\pertanyaan{9}{3. Metode penelitian yang dipilih dan alasan pemilihannya:}
\teksword{8.5}{0}{.........}
\pertanyaan{9}{4. Data, bahan, dan perangkat yang dibutuhkan:}
\teksword{8.5}{0}{.........}
\pertanyaan{9}{5. Gambaran singkat pelaksanaan penelitian yang diusulkan:}
\teksword{8.5}{0}{.........}
\pertanyaan{9}{6. Rencana evaluasi dan indikator keberhasilan:}
\teksword{8.5}{0}{.........}
\teksword{8.5}{0}{.........}

\judulword{10}{4}{Kesimpulan Akhir}
\teksword{8.5}{0}{Rangkum hasil review kedua jurnal: masalah, metode, hasil, kelebihan-keterbatasan, dan alasan usulan pengembangan.}
\teksword{8.5}{0}{.........}
\teksword{8.5}{0}{.........}

\judulword{10}{4}{Daftar Pustaka}
\teksword{8.5}{0}{Cantumkan kedua artikel utama dan referensi pendukung dengan gaya IEEE atau APA secara konsisten.}
\teksword{10}{0}{[1] ..........}
\teksword{10}{0}{[2] ..........}

\judulword{10}{4}{Contoh Format dan Penggunaan}
\pertanyaan{9}{1. Kertas A4; Times New Roman 10–12 pt; spasi 1,0–1,15; margin 2 cm.}
\pertanyaan{9}{2. Contoh panjang laporan: 4 halaman isi; sesuaikan dengan kebutuhan atau ketentuan institusi.}
\pertanyaan{9}{3. Nama file: ReviewJurnal\_NIM\_Nama.pdf.}
\pertanyaan{9}{4. Presentasikan hasil review jika diperlukan sesuai konteks penggunaan.}
\judulword{10}{4}{CONTOH RUBRIK PENILAIAN}
\teksword{8.5}{0}{Komponen dan bobot berikut adalah contoh; sesuaikan dengan kebutuhan, dengan total bobot 100\%.}
\begin{tabelword}{|Y|Y|}
\rowcolor{headerword}
\textbf{\ukuranword{8}Komponen} & \textbf{\ukuranword{8}Bobot} \\ \hline
Analisis masalah kedua jurnal & 15\% \\ \hline
Analisis metode penelitian & 20\% \\ \hline
Analisis data, bahan, dan perangkat & 10\% \\ \hline
Analisis hasil dan bukti & 15\% \\ \hline
Kelebihan dan keterbatasan & 10\% \\ \hline
Perbandingan kritis kedua jurnal & 15\% \\ \hline
Research gap dan usulan pengembangan & 10\% \\ \hline
Kesimpulan, referensi, dan tata tulis & 5\% \\ \hline
TOTAL & 100\% \\ \hline
\end{tabelword}
\teksword{9}{0}{Dosen Pengampu,}
\teksword{9}{24}{(\DosenPengampu)}
\teksword{9}{0}{NIDN: .........}
